% Overleaf-ready section. Dependencies: \usepackage{booktabs,array,hyperref}

\section{Governed Programme-to-Private Follow-Through}
\label{sec:pilot-saved-item-followthrough}

\subsection{Purpose}

Pilot Fabric 0.28 converts a saved television observation into a safe starting
point for personal work.  A product, recipe, destination, song, learning
concept, or idea may be researched, shortlisted, converted into a local draft,
or used to prepare a consequential service proposal.  The architectural rule is:

\begin{quote}
Saving an item records private evidence and interest.  It does not grant
standing authority to purchase, book, message, schedule, or modify an account.
Every consequential continuation requires a new exact scope and a separate
private approval.
\end{quote}

This capability closes the first practical loop from shared television context
to persona-bound private agency while preventing the shared screen from choosing
an identity or external service.

\subsection{Input and Identity Boundary}

Follow-through begins only from an item in state \texttt{saved}.  The private
phone supplies its active \texttt{persona\_id}; the mother reloads the saved item
by \texttt{save\_id} and compares the bound persona before continuing.  A caller
cannot supply an arbitrary title or copy another persona's record into a new
proposal.  Cross-persona research, conversion, approval, refresh, and deletion
are rejected.

The immutable source link consists of:

\begin{verbatim}
{
  "save_id": "save_...",
  "persona_id": "persona_...",
  "category": "product",
  "status": "saved",
  "item_hash": "sha256..."
}
\end{verbatim}

Every continuation record retains the \texttt{save\_id} and source item hash,
making it possible to show which television evidence led to the later research
or proposal without retaining raw programme audio or pixels.

\subsection{Three Continuation Classes}

\subsubsection{Private Research}

Pilot creates a category-specific research query.  Products request evidence on
suitability, price, delivery, returns, safety, and alternatives; destinations
request transport, weather, access, costs, and planning evidence; recipes request
ingredients, timings, substitutions, and food-safety notes.  Music, learning,
and general ideas receive similarly bounded prompts.

The query uses the provider-independent AION research router.  Results may be
served by AION memory, a local model, bounded public evidence, optional Gemini,
or a premium provider according to the configured policy.  Only sanitized HTTPS
links and bounded title, reason, and detail fields enter the private result.
Research sets \texttt{external\_effect} to false and creates no approval.

\subsubsection{Local Drafts}

Shortlists, tasks, and learning plans remain entirely inside the mother brain.
A shortlist records criteria such as suitability, evidence, cost, and risk.  A
task records a title, notes, and incomplete state.  A learning draft records the
topic and source context.  These operations do not contact an account or service
adapter and therefore require no external-action approval.

\subsubsection{Consequential Proposals}

Product, destination, and music saves may prepare shopping, trip, and playlist
proposals.  These use the existing persona-bound Service Execution Hub:

\begin{itemize}
    \item a product creates a shopping proposal with item, quantity, source
    context, and \texttt{purchase: false};
    \item a destination creates a booking proposal with destination, requested
    trip, private date field, and \texttt{booking: false}; and
    \item music creates a music-service proposal with the requested playlist
    concept; the scope explicitly records that no account write occurred.
\end{itemize}

Missing fields place the proposal in \texttt{needs\_details}.  For example, a
trip remains locked until its date is entered on the private phone.  A complete
proposal moves to \texttt{awaiting\_private\_approval}.  The phone displays the
exact parameters and offers explicit Approve and Reject controls.

Approval produces \texttt{approved\_pending\_adapter}, not an execution claim.
An external effect is possible only when the same persona has connected an
authorized adapter, the action is executed once, the uncertain-effect rules pass,
and the provider returns a verified receipt.

\subsection{State Machine}

\begin{center}
\small
Saved item
$\rightarrow$
\begin{tabular}{ll}
research complete & no external effect,\\
local draft ready & no external effect,\\
needs details & private fields missing,\\
awaiting private approval & exact scope complete.
\end{tabular}
\end{center}

The consequential branch continues:

\begin{center}
\small
Awaiting private approval
$\rightarrow$
rejected \quad or \quad approved pending adapter
$\rightarrow$
execute once
$\rightarrow$
verified receipt.
\end{center}

No generative-model result can skip or modify this state machine.

\subsection{Follow-Through Data Contract}

\begin{verbatim}
{
  "schema_version": "pilot.saved-followthrough.v1",
  "followthrough_id": "follow_...",
  "save_id": "save_...",
  "persona_id": "persona_...",
  "kind": "action",
  "action": "shopping",
  "source_save_hash": "sha256...",
  "status": "private_service_proposal_ready",
  "exact_scope": {
    "item": "Extendable ladder",
    "quantity": 1,
    "purchase": false
  },
  "private_approval_created": true,
  "external_effect": false,
  "record_hash": "sha256..."
}
\end{verbatim}

Research records replace \texttt{exact\_scope} with a bounded answer and up to
five sanitized evidence items.  Service records embed the current proposal state
so the private interface and audit can distinguish missing details, pending
approval, rejection, adapter wait, and verified execution.

\subsection{Private Phone Experience}

Each persona-owned saved card now provides four bounded controls:

\begin{enumerate}
    \item \textbf{Research} retrieves evidence without changing an account;
    \item \textbf{Prepare action} selects the category-appropriate shopping,
    trip, playlist, task, or learning route;
    \item \textbf{Make shortlist} creates a local decision aid; and
    \item \textbf{Delete} removes the private saved item.
\end{enumerate}

The \emph{Continue a saved item} panel shows the result route, answer, HTTPS
sources, exact scope, missing private fields, and current approval state.  It
explicitly warns that approval does not equal execution and that no purchase,
booking, message, or external account change is claimed without an adapter
receipt.

\subsection{Security and Privacy Properties}

\begin{itemize}
    \item The shared television cannot launch a persona-bound follow-through.
    \item The source save and active phone persona must match.
    \item Research result URLs must use HTTPS.
    \item Raw provider responses, programme pixels, and raw audio are excluded.
    \item Service parameter validation rejects credential fields including raw
    passwords, tokens, card numbers, CVVs, PINs, and private keys.
    \item Local drafts never acquire an external service binding.
    \item Consequential proposals require a fresh approval unrelated to the
    original Save decision.
    \item Adapter execution is persona-bound, execute-once, and receipt-verified.
    \item Audit records contain identifiers, action class, state, and hashes but
    omit raw private content.
\end{itemize}

\subsection{Verification}

Release 0.28 adds tests for private research, HTTPS sanitization, cross-persona
rejection, shopping-proposal isolation, local draft non-execution, missing trip
details, separate approval, adapter-pending state, and persona-filtered snapshots.
These tests run alongside the complete Fabric suite and the append-only audit-
chain verification.

\subsection{Roadmap Contribution}

This release advances the original master roadmap in two areas.  Under
\emph{Tasks, lists and reminders}, television context can become a private local
task without modifying another person's list.  Under \emph{Shopping, booking
and procurement}, programme evidence can become a private shortlist or exact
proposal while purchase, reservation, payment, and provider execution remain
locked.  Real task-provider connections, baskets, bookings, payments, warranties,
and receipts remain future adapter work rather than claimed functionality.

\subsection{Conclusion}

Programme-to-private follow-through transforms the television from an isolated
answer surface into the beginning of useful personal work.  More importantly,
it proves that this transition can preserve identity, evidence provenance, and
consequential-action boundaries: interest becomes research, research becomes an
exact proposal, and only a new private decision can authorize the next stage.
